\documentclass[a4paper]{article}
\usepackage[fontset=fandol]{ctex}
\usepackage[margin=2.4cm]{geometry}
\usepackage{xcolor}
\usepackage{hyperref}
\usepackage{bookmark}
\usepackage{pdfpages}
\usepackage{tikz}
\usepackage[most]{tcolorbox}
\usetikzlibrary{calc}

\hypersetup{
  unicode=true,
  colorlinks=true,
  linkcolor=blue!50!black,
  urlcolor=blue!50!black,
  pdftitle={CS336 Language Modeling from Scratch 课程讲义合集},
  pdfauthor={Codex},
  pdfsubject={Merged lecture notes with unified cover, page numbers, table of contents, and bookmarks},
  pdfcreator={XeLaTeX},
  bookmarksopen=true,
  bookmarksnumbered=true
}

\setlength{\parindent}{0pt}
\setlength{\parskip}{0.6em}
\setcounter{tocdepth}{1}

\newcommand{\pageoverlay}{%
  \thispagestyle{empty}%
  \begin{tikzpicture}[remember picture,overlay]
    \fill[white] ($(current page.south west)+(0,0)$) rectangle ($(current page.south east)+(0,0.95cm)$);
    \draw[black!20,line width=0.3pt] ($(current page.south west)+(1.8cm,0.95cm)$) -- ($(current page.south east)+(-1.8cm,0.95cm)$);
    \node[anchor=south,font=\small] at ($(current page.south)+(0,0.20cm)$) {\thepage};
  \end{tikzpicture}%
}

\newcommand{\lectureentry}[2]{%
  \pageoverlay
  \phantomsection
  \addcontentsline{toc}{section}{#2}
}

\begin{document}

\begin{titlepage}
\thispagestyle{empty}
\centering
\vspace*{2.6cm}
{\Large Stanford CS336\par}
\vspace{0.4cm}
{\huge\bfseries Language Modeling from Scratch\par}
\vspace{0.5cm}
{\Huge\bfseries 课程讲义合集\par}
\vspace{1.2cm}
{\large Spring 2025\par}
\vspace{1.8cm}

\begin{tcolorbox}[width=0.9\textwidth,colback=blue!4!white,colframe=blue!45!black,sharp corners]
\textbf{内容范围}：第 1 讲至第 17 讲完整讲义合并版\par
\textbf{来源}：\texttt{\detokenize{runs/cs336_all/lectures/}} 下各讲独立 PDF\par
\textbf{合并方式}：统一封面、统一目录、统一页码、PDF 书签\par
\textbf{单讲正文总页数}：174 页\par
\textbf{播放列表}：\href{https://www.youtube.com/playlist?list=PLoROMvodv4rOY23Y0BoGoBGgQ1zmU_MT_}{\nolinkurl{https://www.youtube.com/playlist?list=PLoROMvodv4rOY23Y0BoGoBGgQ1zmU_MT_}}
\end{tcolorbox}

\vfill
{\large 生成时间：\today\par}
\end{titlepage}

\pdfbookmark[0]{目录}{toc}
\setcounter{page}{1}
\tableofcontents
\clearpage

\includepdf[pages={1},pagecommand={\lectureentry{lec01}{CS336 第 1 讲：课程导论、效率世界观与 Tokenization}}]{../lectures/01_overview_and_tokenization/lecture_01_note.pdf}
\includepdf[pages={2-},pagecommand={\pageoverlay}]{../lectures/01_overview_and_tokenization/lecture_01_note.pdf}

\includepdf[pages={1},pagecommand={\lectureentry{lec02}{CS336 第 2 讲：PyTorch、资源核算与训练原语}}]{../lectures/02_pytorch_resource_accounting/lecture_02_note.pdf}
\includepdf[pages={2-},pagecommand={\pageoverlay}]{../lectures/02_pytorch_resource_accounting/lecture_02_note.pdf}

\includepdf[pages={1},pagecommand={\lectureentry{lec03}{CS336 第 3 讲：现代 Transformer 架构与超参数共识}}]{../lectures/03_architectures_hyperparameters/lecture_03_note.pdf}
\includepdf[pages={2-},pagecommand={\pageoverlay}]{../lectures/03_architectures_hyperparameters/lecture_03_note.pdf}

\includepdf[pages={1},pagecommand={\lectureentry{lec04}{CS336 第 4 讲：Mixture of Experts 的机制、代价与工程现实}}]{../lectures/04_mixture_of_experts/lecture_04_note.pdf}
\includepdf[pages={2-},pagecommand={\pageoverlay}]{../lectures/04_mixture_of_experts/lecture_04_note.pdf}

\includepdf[pages={1},pagecommand={\lectureentry{lec05}{CS336 第 5 讲：GPU、内存墙与高性能矩阵乘法直觉}}]{../lectures/05_gpus/lecture_05_note.pdf}
\includepdf[pages={2-},pagecommand={\pageoverlay}]{../lectures/05_gpus/lecture_05_note.pdf}

\includepdf[pages={1},pagecommand={\lectureentry{lec06}{CS336 第 6 讲：Kernel、Triton 与“让 GPU 真正跑快”}}]{../lectures/06_kernels_triton/lecture_06_note.pdf}
\includepdf[pages={2-},pagecommand={\pageoverlay}]{../lectures/06_kernels_triton/lecture_06_note.pdf}

\includepdf[pages={1},pagecommand={\lectureentry{lec07}{CS336 第 7 讲：并行训练基础与 ZeRO / Pipeline / Tensor Parallel}}]{../lectures/07_parallelism_1/lecture_07_note.pdf}
\includepdf[pages={2-},pagecommand={\pageoverlay}]{../lectures/07_parallelism_1/lecture_07_note.pdf}

\includepdf[pages={1},pagecommand={\lectureentry{lec08}{CS336 第 8 讲：跨 GPU 通信与分布式训练原语}}]{../lectures/08_parallelism_2/lecture_08_note.pdf}
\includepdf[pages={2-},pagecommand={\pageoverlay}]{../lectures/08_parallelism_2/lecture_08_note.pdf}

\includepdf[pages={1},pagecommand={\lectureentry{lec09}{CS336 第 9 讲：Scaling Laws I}}]{../lectures/09_scaling_laws_1/lecture_09_note.pdf}
\includepdf[pages={2-},pagecommand={\pageoverlay}]{../lectures/09_scaling_laws_1/lecture_09_note.pdf}

\includepdf[pages={1},pagecommand={\lectureentry{lec10}{CS336 第 10 讲：推理系统}}]{../lectures/10_inference/lecture_10_note.pdf}
\includepdf[pages={2-},pagecommand={\pageoverlay}]{../lectures/10_inference/lecture_10_note.pdf}

\includepdf[pages={1},pagecommand={\lectureentry{lec11}{CS336 第 11 讲：Scaling Laws II}}]{../lectures/11_scaling_laws_2/lecture_11_note.pdf}
\includepdf[pages={2-},pagecommand={\pageoverlay}]{../lectures/11_scaling_laws_2/lecture_11_note.pdf}

\includepdf[pages={1},pagecommand={\lectureentry{lec12}{CS336 第 12 讲：评估体系}}]{../lectures/12_evaluation/lecture_12_note.pdf}
\includepdf[pages={2-},pagecommand={\pageoverlay}]{../lectures/12_evaluation/lecture_12_note.pdf}

\includepdf[pages={1},pagecommand={\lectureentry{lec13}{CS336 第 13 讲：数据工程 I}}]{../lectures/13_data_1/lecture_13_note.pdf}
\includepdf[pages={2-},pagecommand={\pageoverlay}]{../lectures/13_data_1/lecture_13_note.pdf}

\includepdf[pages={1},pagecommand={\lectureentry{lec14}{CS336 第 14 讲：数据工程 II}}]{../lectures/14_data_2/lecture_14_note.pdf}
\includepdf[pages={2-},pagecommand={\pageoverlay}]{../lectures/14_data_2/lecture_14_note.pdf}

\includepdf[pages={1},pagecommand={\lectureentry{lec15}{CS336 第 15 讲：对齐：SFT / RLHF}}]{../lectures/15_alignment_sft_rlhf/lecture_15_note.pdf}
\includepdf[pages={2-},pagecommand={\pageoverlay}]{../lectures/15_alignment_sft_rlhf/lecture_15_note.pdf}

\includepdf[pages={1},pagecommand={\lectureentry{lec16}{CS336 第 16 讲：对齐：RL I}}]{../lectures/16_alignment_rl_1/lecture_16_note.pdf}
\includepdf[pages={2-},pagecommand={\pageoverlay}]{../lectures/16_alignment_rl_1/lecture_16_note.pdf}

\includepdf[pages={1},pagecommand={\lectureentry{lec17}{CS336 第 17 讲：对齐：RL II}}]{../lectures/17_alignment_rl_2/lecture_17_note.pdf}
\includepdf[pages={2-},pagecommand={\pageoverlay}]{../lectures/17_alignment_rl_2/lecture_17_note.pdf}

\end{document}
